\section{Topological classification and semantic context}
\label{sec:model-theory}
The four degree types have a purely topological description valid for
all the uniform requests of Theorem~\ref{thm:classification}. It uses
only the eventual prime-power periodicity of bounded-degree
sublevel sets. No Boolean encoding or model-existence assumption is
needed. We then recall the distinction between algebraic witnesses
and actual solutions, and the established Euler-characteristic criterion.

\subsection{The \texorpdfstring{$p$}{p}-adic hardness spectrum}
\label{sec:topological-spectrum}
The underlying structure is already implicit in the residue-class
analysis of my 2009 report~\cite[Section~4, pp.~18--19]{RiisLong}:
once a degree threshold is fixed, sufficiently large sizes are classified
by residues modulo a power of $p$. We express that earlier structure
in $p$-adic language and apply it here to uniform consequences as well
as refutations.

Fix any one of the aggregate or fixed-orbit profiles $h$ in
Theorem~\ref{thm:periodicity}, over a field of characteristic $p$.
A $p$-adic integer is a compatible choice of residues modulo all powers
of $p$. The space $\mathbb Z_p$ is compact, and its basic residue
neighbourhoods $a+p^L\mathbb Z_p$ are both open and closed
(\emph{clopen}). Using the onsets $N(d)$ and periods $Q(d)$ from
\eqref{eq:NQ}, let $A_d(h)\subseteq\{0,\ldots,Q(d)-1\}$ consist of
the residues $a$ for which
\[
 h(n)\le d\quad\text{whenever }n\ge N(d),\qquad
                         n\equiv a\pmod{Q(d)}.
\]
Put
\[
 \mathcal E_d(h)=\bigcup_{a\in A_d(h)}
                         \bigl(a+Q(d)\mathbb Z_p\bigr).
\]
Thus $\mathcal E_d(h)$ records the eventually degree-$d$-easy residue
classes. Each is clopen, and $\mathcal E_d(h)\subseteq\mathcal E_e(h)$
for $d\le e$: refine to the common period and use the periodicity on a
sufficiently late tail. Set
\begin{equation}\label{eq:hardness-spectrum}
 \mathcal H_p(h)=\mathbb Z_p\setminus
                    \bigcup_{d\in\mathbb N}\mathcal E_d(h).
\end{equation}
We call this closed set the \emph{hardness spectrum}. It does not depend
on the choice of valid onsets or prime-power periods. Equivalently,
$x\notin\mathcal H_p(h)$ exactly when some residue neighbourhood of
$x$ has a uniform finite degree bound at all sufficiently large integer
sizes in that neighbourhood. Common residue refinements show the
agreement of this description with \eqref{eq:hardness-spectrum}.

\begin{samepage}
\begin{corollary}[The four topological shapes]
\label{cor:topological-spectrum}
For every uniform request considered in this paper, the NS and PC
profiles have the following classification:
\begin{center}
\begin{tabular}{@{}ll@{}}
\toprule
Type & Hardness spectrum $\mathcal H_p(h)$\\
\midrule
$\mathrm T_1$ & Empty\\
$\mathrm T_2$ & All of $\mathbb Z_p$\\
$\mathrm T_3$ & Nonempty, proper and clopen\\
$\mathrm T_4$ & Closed but not open\\
\bottomrule
\end{tabular}
\end{center}
Moreover, for the same request the two systems have the same hardness
spectrum, although their exact degrees and threshold sets can differ.
\end{corollary}
\end{samepage}
\begin{proof}
Write $E_d=\mathcal E_d(h)$ and $H=\mathcal H_p(h)$. If $H$ is empty,
compactness and the increasing open cover $(E_d)_d$ give
$E_D=\mathbb Z_p$ for some $D$. This is eventual boundedness, or T1;
the converse is immediate. Also $H=\mathbb Z_p$ exactly when every
$E_d$ is empty. This says that no finite degree recurs, or T2.

In T3, let $D$ be the largest recurring finite degree.
Theorem~\ref{thm:classification} gives a fixed nonempty proper union
of easy residue classes. It is $E_D$, and no $E_d$ meets a hard class:
otherwise arbitrarily late integers in a common residue refinement
would have both bounded and linearly growing degree. Hence
$\bigcup_dE_d=E_D$, so $H$ is nonempty, proper and clopen.

Conversely, suppose $H$ is nonempty, proper and clopen. Its compact
complement is covered by the increasing $E_d$, hence equals some $E_D$.
If a finite degree $e$ recurs, periodicity gives a residue neighbourhood
on which that exact value occurs on a tail. This neighbourhood lies
in $E_e\subseteq E_D$; comparing late integer representatives gives
$e\le D$. There is a recurring finite degree because the complement
of $H$ is nonempty, but $H\ne\varnothing$ excludes T1. This is T3.
Since $H$ is always closed and the four degree types are exhaustive,
the remaining shape, closed but not open, is exactly T4.

Finally, the comparison theorem gives
$h_{\PC}\le h_{\NS}\le F(h_{\PC})$, with $F$ independent of size.
Thus a residue neighbourhood is eventually bounded in one system if
and only if it is eventually bounded in the other. The equivalent
description above proves equality of their hardness spectra.
\end{proof}

In T3 a single finite congruence resolution separates the easy and
hard regions. In T4 the closed spectrum has a nonempty boundary:
every neighbourhood of a boundary point meets the easy region, but
no such neighbourhood has one uniform finite degree bound.
The spectrum records the qualitative organisation of the profile;
the linear lower bounds and finite linear peaks still come from the
quantitative transfer theorem. It is defined from the polynomial
request and its degree profile, not from the ordinary model class
of a first-order sentence alone.

\subsection{Dual witnesses for failure to derive a target}
A linear functional is a scalar-valued linear map on coefficient
vectors. If it vanishes on every permitted consequence but not on
the target, applying it to a proposed derivation gives a contradiction.
The converse is finite-dimensional separation by a quotient basis;
it does not require the functional to be evaluation at a solution.

For either system write $W_d(n)$ for its degree-$d$ consequence space:
$W_d(n)=C_d(n)$ for PC and $W_d(n)=I_d(n)$ for NS. Let
$\delta=D_G$ be the fixed upper bound on target degree. We view
$W_d(n)$ inside $P_{\le\max\{d,\delta\}}(n)$, so the ambient space
also contains every requested target even when its degree exceeds $d$.

\begin{definition}[A degree-bounded target witness]
\label{def:consistency-witness}
For a requested target $G_{b,n}$, a degree-$d$ witness is a $K$-linear
functional
\[
 \lambda:P_{\le\max\{d,D_G\}}(n)\longrightarrow K,\qquad
 \lambda(G_{b,n})=1,\qquad\lambda(W_d(n))=0.
\]
For an all-target request, a witness consists of an index $b$ and such
a functional. For a single fixed orbit, one may use any available
representative. With target $1$, this reduces to the usual normalisation
$\lambda(1)=1$ for a bounded-degree non-refutability witness.
\end{definition}

\begin{proposition}[Dual form of nonderivability]
\label{prop:dual-witness}
A degree-$d$ target witness exists exactly when the target has no
derivation of degree at most $d$ in the chosen system. For an all-target
request it exists exactly when $h(n)>d$.
\end{proposition}
\begin{proof}
If $G_{b,n}\in W_d(n)$, a functional annihilating $W_d(n)$ cannot
take value one on it. Conversely its class in the finite-dimensional
quotient $P_{\le\max\{d,D_G\}}(n)/W_d(n)$ is nonzero whenever
it is not in $W_d(n)$. Extend that class to a basis, assign it value
one and assign the other basis vectors zero. Composing with the
quotient map gives the required functional. If the target degree is
greater than $d$, it automatically lies outside $W_d(n)$; the enlarged
ambient space also covers this degree-floor obstruction. Failure of
all-target derivability is failure for at least one target.
\end{proof}

A $K$-valued solution $a$ of $\Gamma_n$ which violates $G_{b,n}=0$ gives a
witness by $\lambda(f)=f(a)/G_{b,n}(a)$. For target one this is ordinary
evaluation, with its multiplicative property. A general witness need
not be multiplicative, positive or normalised on one. It records a
linear separation from all permitted bounded-degree consequences,
not necessarily a truth assignment to the original variables.

\begin{example}[Semantic entailment need not be derivation]
The equation $x^2=0$ over a field forces $x=0$ at every solution.
Nevertheless, neither NS nor PC derives $x$ from this axiom: every
consequence belongs to the ideal $(x^2)$. For $d\ge2$, the functional
on $P_{\le d}$ taking value one on $x$ and zero on all other monomials
annihilates the degree-$d$ consequence space. This is a target witness
although no actual solution violates $x=0$. The example explains why
our minimum degree is a syntactic ideal-membership resource.
\end{example}

\begin{example}[Linear consistency without truth-value evaluation]
Over $\mathbb F_3$, take $x^2-x=0$, target one and degree budget two.
The PC consequence space is $K(x^2-x)$: no nonzero vector in that span
has degree at most one, so no multiplication is usable. The functional
\[
 \lambda(1)=1,\qquad\lambda(x)=\lambda(x^2)=2
\]
is a witness but is not evaluation: $\lambda(x^2)=2$ whereas
$\lambda(x)^2=1$. Boolean equations impose the intended truth-value
semantics on actual solutions, without making every linear witness an
actual evaluation.
\end{example}

For target $1$, the NS witness is the design method for lower bounds
surveyed and generalised by Buss~\cite{BussDesign}: a normalised
linear functional annihilates every permitted axiom multiple.
Annihilating only axiom multiples gives an NS witness. For PC the
functional must annihilate the full closure, including polynomials
obtained by cancellation followed by new multiplication. This is the
dual form of the distinction between $I_d$ and $C_d$.

\subsection{What Kraj\'i\v cek's Euler criterion supplies}
A strong Euler characteristic assigns a value $\chi(A)$ to each
permitted parameter-definable set, possibly infinite. It gives a point
value one, is additive on disjoint unions, is multiplicative on products,
and is invariant under definable bijections. For a definable map
$f:A\to B$ whose fibres all have characteristic $r$, it also satisfies
$\chi(A)=r\chi(B)$. This last fibre rule is essential.

Such a rule replaces finite counting in a definable solution of a
uniform polynomial system. In Kraj\'i\v cek's criterion, the values lie
in $\mathbb Z/p^\nu\mathbb Z$, where the required precision $\nu$
depends on the degree budget and the indexing bounds. Reduction modulo
$p$ interprets a sum of a definable $\mathbb F_p$-valued function by
\[
 \sum_{x\in A}^{\chi}a(x)
       =\sum_{b\in\mathbb F_p}b\,\overline\chi(a^{-1}(b)).
\]
With the definability and index-sort hypotheses of
\cite[Theorem~6.1]{KrajEuler}, including the characteristic on the
required index sorts $M^*$, a solution violating the target $g=0$
excludes derivations within that fixed degree budget at sufficiently
large sizes in the residue class specified by $\chi(M)\bmod p^\nu$.
The refutation case is $g=1$.

The mechanism is that bounded-degree calculations involve only bounded
configurations of indices. The counting rules transport the relevant
finite identities to the definable setting. A hypothetical derivation
would make the target vanish there, contradicting the chosen solution.
A characteristic with sufficient precision for each degree therefore
provides a semantic obstruction to bounded-degree proofs. As the degree
increases, the required precision may increase too.

Definable bijections, disjoint-union relations and products can be
organised in the Grothendieck ring of definable sets
\cite{KrajScanlon}. The strong fibre rule is additional structure.
Neither an arbitrary infinite model nor a merely mod-$p$ counting rule
supplies the full hypotheses of Kraj\'i\v cek's criterion. It is a
sufficient criterion; no converse is used or asserted here.

Richer first-order counting theories can also describe bounded-degree
algebraic reasoning under additional hypotheses~\cite{PTT}. Such
model-theoretic representations are not needed for
Corollary~\ref{cor:topological-spectrum}; in particular, its scope is
not restricted to a fixed Boolean compiler. A general converse to the
Euler criterion remains a separate question.
